\chapter{The Journal-Relative Web}
\label{sec:journal-web}
\label{sec:paths}
\label{sec:policy}

\section{Resources and Rooted Paths}

A resource is a leaf reached through path and object semantics. Its content may be uninterpreted bytes, an encoded expression, or another self-encoding structure. Resource identity has two complementary dimensions. A path gives the resource continuity within one selected journal-relative coordinate system, while a digest identifies the particular state resolved at that path.

Interpretation always begins from a root. A path without its root is incomplete in the same way that a relative URL is incomplete without a base. Traversal follows documented semantic boundaries rather than assuming one universal physical tree, so two objects may interpret the same path segment differently while still preserving an end-to-end integrity proof. Resolution itself proves neither existence, availability, principal identity, nor authority until the relevant structural and policy checks succeed.

\section{Plural Coordinates}

No global top-level registry is required. This follows traditions of linked local namespaces and self-certifying pathnames~\cite{Rivest1996sdsi,Mazieres1998path}. Each journal induces coordinates over the local and foreign state reachable from its root. Identical suffixes beneath two roots need not name the same resource, and a bridge segment records the relationship through which foreign state became reachable. A historical index additionally selects the committed head against which the rest of the path is resolved.

Traversal can end at a stub, an absent resource, an unreachable peer, private terminal state, or a policy denial. These outcomes have different meanings and should remain distinguishable to clients. Content identity, path continuity, human-readable naming, and principal identity are likewise separate concerns. Root-relative addressing complements content-bound and data-centric naming rather than replacing them~\cite{Haber1997secure,Kuhn2014trusty,Zhang2014named}.

Eliminating one mandatory naming authority moves work into root selection, discovery, route maintenance, and comparison among viewpoints. The network shown in \cref{fig:journal-relative-web} is therefore journal-relative: independently advancing histories are joined by sparse explicit relationships, with no enclosing global root from which every resource must descend.

\section{Web Projection}

Addressable network form and query-optimized operational form serve different purposes. The first favors independently identifiable resources, linking, selective transfer, citation, and plural ownership. The second favors application records, schemas, indices, transactions, compression, and local query performance. Neither is a universally superior physical representation; \cref{tab:data-forms} summarizes their complementary roles.

\begin{table}[htbp]
  \centering
  \small
  \begin{tabularx}{\textwidth}{@{}p{0.18\textwidth}XX@{}}
\toprule
\textbf{Aspect} & \textbf{Addressable Network Form} & \textbf{Query-Optimized Operational Form} \\
\midrule
Primary unit & Independently identified resource and relationship & Application record, row, object, index, or transaction \\
Identity & Root-relative path plus resolved state commitment & Application key interpreted within one operational system \\
Strengths & Linking, citation, selective transfer, plural ownership, and historical reference & Locality, transactions, compression, indexing, and efficient bounded queries \\
History & Durable resource states and cross-history relationships are first-class & Commonly application-defined, compacted, or maintained in separate logs \\
Typical costs & Traversal, synchronization, repeated metadata, and derived indexing & Provider dependence, schema coupling, and application-specific portability \\
Architectural role & Shared network projection and durable evidence & Load-bearing operations and rebuildable derived views \\
\bottomrule
\end{tabularx}

  \caption{Complementary Data Forms}
  \label{tab:data-forms}
\end{table}

A Web projection lets an operational system remain load-bearing while selected state receives a durable addressable form. The projection must define resource boundaries, deterministic encodings, path conventions, progress manifests, and the semantics of correction, privacy, deletion, and replay. It does not automatically inherit every invariant of its source. Transaction boundaries, relational constraints, and mutable deletion policies may require explicit representation in the projected history.

Adoption can therefore proceed incrementally. An asynchronous projection first supplies an audit path. As clients, authorization, query facilities, and monitoring mature, selected projected resources may become load-bearing. Observable cursors and deterministic correction records make lag and divergence explicit, but coexistence still costs synchronization, reconciliation, and duplicated storage. \WIP General projection contracts for dynamic applications, privacy-sensitive records, correction, and bidirectional synchronization remain unsettled.

\section{Historical Identity and Keys}

\WIP A remote principal becomes meaningful only when the receiving journal can resolve the origin through its own bridge graph to a current public interface key. A synchronized descriptor binds the remote journal's name, stable root key, current interface key, canonical endpoints, head, and historical context. Unverified information endpoints and private configuration may cache those facts for convenience, but they cannot override the signed state.

The root key signs journal heads and is a pinned, non-rotating trust anchor in the initial design. Interface keys sign federated invocations and may change after their replacements appear in root-signed heads synchronized through the bridge graph. A compiled return proof reveals the exact origin checkpoint and interface key that the terminal can verify. The origin can use that retained key or complete its final local proof segment to a newer key authorized by a later root-signed head.

A stored head provides reverse cryptographic reachability even when its journal cannot accept inbound transport. Cryptographic reachability, network reachability, and responsibility for initiating synchronization are therefore distinct. Reciprocal heads, descriptors, continuity evidence, and route proofs form a public control plane; none grants authority over private application state.

\section{Federated Route Establishment}

\WIP A private federated invocation begins with one public multi-hop request and response that compile two different proofs. Each proof is rooted in the current state of the endpoint that will consume it. For an origin A and terminal B, the forward leg constructs A's proof of B. A begins with the exact adjacent-journal checkpoint stored in A's current head. Every intermediary proves continuity from the checkpoint it was actually given to its current head, then contributes the next nested checkpoint stored there. B completes the proof to its current root identity, interface key, head, descriptor, and canonical endpoint.

The return leg starts from B's own current head and follows the reverse checkpoints in the same manner. When the response reaches A, A completes the final segment to the interface key it will use to sign. This endpoint-relative construction prevents an upstream journal from guessing what a later participant has synchronized. Every link consists of committed Merkle inclusion and signed chain continuity; no live intermediary attestation is treated as though it had already been committed into the terminal's graph.

The resulting B-relative source proof is mandatory on the direct invocation. Beginning from its own current head, B can verify the proof, resolve the exact A interface key selected by it, verify A's signature, construct the receiver-relative principal, authorize the function, and execute it in one pass. This realizes a more abstract verifiable-chain idea: identity is established through a connected path that the relying endpoint can verify from its own state and through which evidence is traced in both directions~\cite{dinh2025identity}.

After the public handshake, A sends the signed function, arguments, private paths, and values directly to B over HTTPS. Intermediate journals handle only public proof and discovery material; they neither relay nor observe the private call. The design consequently requires the terminal to be reachable and deliberately excludes onion routing, encrypted offline relay, and offline delivery of private invocations. The origin signature binds the complete semantic route, terminal audience, explicit identity, function, and arguments. Local credentials and private signing keys remain inside the origin journal boundary.

\Cref{fig:federated-invocation} separates the compiled public route from the direct final mile. The self-contained proof removes the need for a terminal route cache, compact terminal-issued receipt, or second reverse traversal. \WIP Exact proof serialization, bridge-alias translation, endpoint and head races, retry behavior, and terminal response binding remain to be specified.

\begin{figure}[htbp]
  \centering
  \resizebox{0.88\textwidth}{!}{% Public routing follows exact nested journal objects already committed beneath
% the origin head.  The private invocation then bypasses the intermediary.
\begin{tikzpicture}[
  participant/.style={draw=swblue, thick, rounded corners=3pt, fill=swlightblue,
                      minimum width=31mm, minimum height=8mm,
                      font=\small\bfseries, align=center},
  lifeline/.style={draw=swgray!65, densely dotted},
  msg/.style={-{Latex[length=2.2mm]}, thick, draw=swgray},
  return/.style={msg, dashed},
  direct/.style={msg, draw=swblue, very thick},
  note/.style={font=\footnotesize, text=swgray!85!black, align=center},
]
  \node[participant] (a) at (-5.0,0) {Origin A};
  \node[participant] (c) at (0,0) {Intermediary C};
  \node[participant] (b) at (5.0,0) {Terminal B};
  \draw[lifeline] (a.south) -- (-5.0,-7.0);
  \draw[lifeline] (c.south) -- (0,-4.9);
  \draw[lifeline] (b.south) -- (5.0,-7.0);

  \draw[msg] (-5.0,-1.0) -- node[above,note]
    {follow exact C object in selected A head} (0,-1.0);
  \draw[msg] (0,-1.8) -- node[above,note]
    {follow exact nested B object; preserve digest} (5.0,-1.8);

  \draw[return] (5.0,-2.8) -- node[above,note]
    {B-rooted partial object + reverse route to A key} (0,-2.8);
  \draw[return] (0,-3.6) -- node[above,note]
    {same B index/digest + root-sequence context} (-5.0,-3.6);
  \node[note, anchor=east] at (-5.1,-4.15)
    {A matches returned B object to its committed forward route};

  \draw[swgray!55, dashed] (-6.3,-4.8) -- (6.3,-4.8);
  \node[note, fill=white, inner sep=2pt] at (0,-4.8)
    {direct final mile; intermediary sees no private function, path, arguments, or value};

  \draw[direct] (-5.0,-5.65) -- node[above,note, text=swblue]
    {B-rooted object + signed get/set!/resolve} (5.0,-5.65);
  \draw[return, draw=swblue] (5.0,-6.55) -- node[above,note, text=swblue]
    {live get/set! result; resolve returns verified proof} (-5.0,-6.55);
\end{tikzpicture}
}
  \caption{Compiled Route Proof and Invocation}
  \label{fig:federated-invocation}
\end{figure}

\section{Journal Boundary and Authorization}

External requests enter through a controlled journal interface responsible for authentication, authorization, argument normalization, route preparation, batching, and explicit persistence. Network fetches and service credentials remain outside deterministic object methods. Administration over code, roots, bridges, secrets, and policy remains local even when ordinary resource operations become federated.

\WIP Invocation identity is explicit and scalar. It identifies the originating journal, one user locally authenticated by that journal, or a public caller. Omission is invalid and can never imply privileged journal identity. At the receiver, the verified route and scalar identity combine into the canonical receiver-relative principal path used by authorization.

When the scalar denotes a local user, the origin attests that it authenticated that user through its own session or token system. The terminal does not validate the remote password or session. It first authenticates the origin journal and proof, then decides what authority to grant the narrower user attestation. Every terminal operation and every dynamically discovered read requires an explicit local rule; bridge reachability, signatures, and structural integrity are not permission.

Prefix rules express ordinary ownership efficiently, but they may not capture object methods, callbacks, or effects spanning several paths. Principal authentication should also remain distinct from verifiable-credential infrastructure~\cite{Sedlmeir2021digital}. \WIP The signed sender flow, authenticated response envelope, arbitrary object methods, links, and programmable reads still require end-to-end inspection.

\section{Replay and Mutation}

Federation does not promise exactly-once execution. It maintains no generic durable request-ID table, nonce protocol, timestamp window, or replay cache, so a valid signed invocation may be replayed. Reads, route proofs, and head synchronization must therefore be idempotent, with ordinary rate and resource controls limiting repeated expensive work.

A mutation must either be naturally replay-safe, explicitly permit repetition, or carry a signed operation-specific precondition such as an expected revision or state digest. A successful guarded mutation advances that state, so replaying the old signed request fails its precondition rather than restoring stale data. Operations with irreversible external effects need stronger domain-specific idempotency and should not become federated merely because their request can be signed. \WIP The first federated mutation set and the precondition semantics for each function remain to be selected.

\section{Links and Discovery}

\WIP A first-class link identifies a target through an explicit root and relationship context and states whether resolution is current or historical. The integrity of the link and the integrity of the target are separate. Traversal must preserve route and authority boundaries and distinguish a broken link from private, unauthorized, malformed, unavailable, or internally inconsistent state. Unrestricted traversal also creates privacy, denial-of-service, and confused-deputy risks.

Discovery finds candidate journals, resources, endpoints, and routes through social exchange, indices, descriptors, existing bridges, or out-of-band knowledge. A discovered candidate does not thereby become accepted or authoritative. Candidates may be stale, mutually inconsistent, or malicious; local policy decides which relationships to create and which routes to follow. Endpoint and interface-key changes require root-signed historical continuity because discovery metadata evolves. Root-key recovery and rotation remain deferred.

\section{Transparent Clients and Adapters}

The primary interaction model is ordinary data rather than proof ceremony. Clients read and write resources, invoke objects, follow links, and resolve history. A specialized local client chooses roots and historical context, fetches and merges evidence, enforces policy, reports meaningful failure distinctions, and applies retention intent beneath those operations.

Cryptography narrows dependence on remote operators but does not eliminate the local trusted computing base. Users still depend on their client, evaluator, operating system, and key store. Gateways, filesystems, browsers, inspection tools, and projection workers are replaceable adapters over common resource and history semantics, yet any adapter can hide richer behavior or misrepresent a result and therefore remains within a client or service trust boundary.

Application protocols and provenance vocabularies belong above the generic structures~\cite{Moreau2010open,Groth2010anatomy,BernersLee2001semantic}. \WIP Native browsers and computational views should make routine verification unobtrusive while preserving explicit audit and debugging views for specialists.
